\section[Xây dựng chương trình]{Xây dựng chương trình (Dashboard, bài học)}

\subsection{4 Xây dựng chương trình}

\begin{frame}{Dashboard: bốn dashboard, bốn loại}
\muc{4 · Xây dựng chương trình}
\begin{itemize}
  \item DB1 Tổng quan và doanh thu
  \item DB2 Giao hàng và trải nghiệm khách
  \item DB3 Địa lý và phí vận chuyển
  \item DB4 Thu hút và hiệu quả seller
\end{itemize}
\khungchen[0.3\textheight]{CHÈN HÌNH / BẢNG: bảng 4 dashboard (\texttt{huong-dan/dashboards/})}
\note{[S89] Mỗi dashboard dùng nhiều loại biểu đồ, slicer và định dạng có điều kiện. Tiêu đề đầy đủ trước khi rút gọn: Bốn dashboard theo bốn loại: dòng thời gian có dự báo, thống kê, OLAP và vận hành Tiêu đề trước đây: Bốn dashboard, bốn loại: dòng thời gian có dự báo, thống kê, OLAP, vận hành}
\end{frame}

\begin{frame}{Mô hình Power BI: 16 bảng sao, 32 quan hệ cho DB1–DB4}
\muc{4 · Xây dựng chương trình}
\begin{itemize}
  \item Quan hệ còn lại không hoạt động
\end{itemize}
\khungchen[0.3\textheight]{CHÈN HÌNH / BẢNG: ảnh chụp Model view}

\note{[S90] Tham số \texttt{DataFolder} giúp đổi máy chỉ cần sửa một giá trị. Tiêu đề trước đây: Mô hình Power BI có 16 bảng và 32 quan hệ, mỗi fact một quan hệ ngày hoạt động}
\end{frame}

\begin{frame}{Nguyên tắc thiết kế: bố cục và định dạng}
\muc{4 · Xây dựng chương trình}
\begin{itemize}
  \item Mỗi trang có thanh tiêu đề, icon minh hoạ và menu điều hướng giữa các trang
  \item Mỗi đồ thị có tiêu đề
  \item Màu thống nhất: xanh là tốt, vàng cần chú ý, đỏ là xấu
\end{itemize}
\note{[P12] Checklist của thầy: màu sắc bố cục hài hoà, tiêu đề dashboard và đồ thị, icon minh hoạ, menu. Theme ở powerbi/theme.json; ngưỡng cảnh báo on-time 90\%/95\%.}
\end{frame}

\begin{frame}{Nguyên tắc thiết kế: chủ đề màu}
\muc{4 · Xây dựng chương trình}
\begin{itemize}
  \item \texttt{powerbi/theme.json}; ngưỡng cảnh báo on-time 90\%/95\%
\end{itemize}
\khungchen[0.3\textheight]{CHÈN HÌNH / BẢNG: ảnh chụp theme}

\note{[S92] Cùng một quy ước màu trên mọi dashboard. Tiêu đề trước đây: Chủ đề màu thống nhất: xanh là tốt, vàng cần chú ý, đỏ là xấu}
\end{frame}

\begin{frame}{Quy chuẩn định dạng số: con số có nghĩa}
\muc{4 · Xây dựng chương trình}
\begin{BangTC}{lXr}
\textbf{Loại số} & \textbf{Quy tắc} & \textbf{Ví dụ} \\ \midrule
Tiền & rút gọn theo triệu hoặc nghìn & 13,5 triệu \\
Số đơn & rút gọn theo nghìn & 98 nghìn \\
Phần trăm trên thẻ KPI & làm tròn số nguyên & 93\% \\
Phần trăm dưới 10\% hoặc cần phân biệt & một chữ số thập phân & 3,1\% \\
\end{BangTC}
\KetLuan{Chỉ giữ chữ số có nghĩa: 12,39\% hiển thị 12\%}
\note{[P12] Checklist của thầy: đơn vị tính trên dashboard, định dạng số, dùng con số có nghĩa (ví dụ 12,39\% thành 12\%). Giá trị thật vẫn đầy đủ trong DAX; đây chỉ là cách hiển thị. Quy chuẩn ghi ở huong-dan/dashboards/README.md.}
\end{frame}

\begin{frame}{DB1 – Tình huống: ban điều hành hỏi gì}
\muc{4 · DB1 Tổng quan và doanh thu}
\begin{itemize}
  \item Ai dùng: ban điều hành và quản lý kinh doanh, xem hằng tuần và hằng tháng
  \item Câu hỏi: GMV tăng trưởng MoM, YoY và quy luật 80/20
\end{itemize}
\note{[P12] Ai dùng và câu hỏi của DB1 (mạch Tình huống, Dashboard, Phát hiện, Quyết định).}
\end{frame}

\begin{frame}{DB1 – Dashboard: trang 1 tổng quan}
\muc{4 · DB1 Tổng quan và doanh thu}
\begin{itemize}
  \item Thẻ KPI, đường GMV theo tháng, cột MoM, thanh theo vùng, treemap ngành
\end{itemize}
\khungchen{CHÈN ẢNH CHỤP DB1 trang 1}

\note{[S93] Trang tổng quan cho lãnh đạo. Tiêu đề trước đây: DB1 trang 1: GMV 13,49 triệu BRL, 98.199 đơn, đúng hạn 93,2\%, review 4,09}
\end{frame}

\begin{frame}{DB1 – Dashboard: trang 1 xu hướng và dự báo}
\muc{4 · DB1 Tổng quan và doanh thu}
\begin{itemize}
  \item Black Friday 152.654 BRL trong một ngày
\end{itemize}
\khungchen{CHÈN ẢNH CHỤP DB1 trang 1}

\note{[S94] Từ 03/2018 GMV đi ngang nhưng cao hơn cùng kỳ (05/2018 +97\%). Tiêu đề đầy đủ trước khi rút gọn: DB1 trang 1: tháng 11/2017 là đỉnh 1,0 triệu BRL, Sudeste 65\% GMV, có đường dự báo Tiêu đề trước đây: DB1 trang 1: đỉnh 11/2017 là 1,0 triệu BRL, Sudeste 65\% GMV, có dự báo}
\end{frame}

\begin{frame}{DB1 – Dashboard: trang 2 danh mục}
\muc{4 · DB1 Tổng quan và doanh thu}
\begin{itemize}
  \item Treemap, ma trận drill, Pareto, scatter AOV – freight
\end{itemize}
\khungchen{CHÈN ẢNH CHỤP DB1 trang 2}

\note{[S95] Drill từ nhóm ngành xuống danh mục. Tiêu đề trước đây: DB1 trang 2: 18 danh mục tạo 80\% GMV; Home \& Furniture dẫn đầu 22,8\%}
\end{frame}

\begin{frame}{DB1 – Dashboard: trang 2 phí vận chuyển}
\muc{4 · DB1 Tổng quan và doanh thu}
\begin{itemize}
  \item Freight ratio 16,6\% toàn sàn
\end{itemize}
\khungchen[0.3\textheight]{CHÈN HÌNH / BẢNG: ảnh chụp DB1 trang 2 (scatter)}

\note{[S96] Danh mục nhẹ giá thấp bị phí vận chuyển ăn vào nhiều nhất. Tiêu đề trước đây: DB1 trang 2: phí vận chuyển nặng ở hàng giá thấp và ở vùng xa}
\end{frame}

\begin{frame}{DB1 – Phát hiện: ba số cần nhớ}
\muc{4 · DB1 Tổng quan và doanh thu}
\begin{itemize}
  \item GMV 13,49 triệu BRL, 98.199 đơn, đúng hạn 93,2\%, review 4,09
  \item Tháng 11/2017 là đỉnh 1,0 triệu BRL; Sudeste 65\% GMV; Black Friday 152.654 BRL trong một ngày
  \item 18 danh mục tạo 80\% GMV, Home \& Furniture dẫn đầu 22,8\%
\end{itemize}
\note{[P12] Ba phát hiện có số của DB1; số lấy từ các slide Dashboard cùng nhóm và từ phần khám phá dữ liệu.}
\end{frame}

\begin{frame}{DB1 – Quyết định: hành động đề xuất}
\muc{4 · DB1 Tổng quan và doanh thu}
\begin{itemize}
  \item Chuẩn bị tồn kho và vận chuyển trước Black Friday
  \item Tập trung vào 18 danh mục tạo 80\% GMV
  \item Xem lại phí vận chuyển ở hàng giá thấp (freight ratio 16,6\%)
\end{itemize}
\note{[P12] Đề xuất hành động rút từ ba phát hiện; nhóm điều chỉnh theo cách hiểu nghiệp vụ của mình.}
\end{frame}

\begin{frame}{DB1 – DAX chính}
\muc{4 · DB1 Tổng quan và doanh thu}
\khungchen[0.3\textheight]{CHÈN ẢNH / ĐOẠN DAX: measure chính của DB1 (powerbi/measures.md)}
\note{[P12] Khung chờ ảnh đoạn DAX chính; nguồn powerbi/measures.md (Người phụ trách dashboard điền).}
\end{frame}

\begin{frame}{DB2 – Tình huống: vận hành và chăm sóc khách}
\muc{4 · DB2 Giao hàng và trải nghiệm khách}
\begin{itemize}
  \item Ai dùng: vận hành và logistics (hằng ngày đến hằng tuần), chăm sóc khách hàng
  \item Câu hỏi: hứa so với thực tế chênh bao nhiêu, giao trễ có kéo điểm review xuống không
\end{itemize}
\note{[P12] Ai dùng và câu hỏi của DB2 (mạch Tình huống, Dashboard, Phát hiện, Quyết định).}
\end{frame}

\begin{frame}{DB2 – Dashboard: trang 1 lead time và phễu}
\muc{4 · DB2 Giao hàng và trải nghiệm khách}
\begin{itemize}
  \item Funnel, cột chồng từng chặng, bản đồ trễ
\end{itemize}
\khungchen{CHÈN ẢNH CHỤP DB2 trang 1}

\note{[S97] KPI "hứa so với thực tế" nằm trên đầu trang. Tiêu đề trước đây: DB2 trang 1: lead time 12,6 ngày, phễu trạng thái và tỷ lệ trễ theo bang}
\end{frame}

\begin{frame}{DB2 – Dashboard: trang 1 hứa so với thực tế}
\muc{4 · DB2 Giao hàng và trải nghiệm khách}
\begin{itemize}
  \item Đúng hạn tụt 87,6\% ở Black Friday 11/2017 và 81,0\% ở 03/2018
\end{itemize}
\khungchen{CHÈN ẢNH CHỤP DB2 trang 1}

\note{[S98] Tải cao làm giao chậm dù trung bình vẫn "đúng hạn". Tiêu đề trước đây: DB2 trang 1: Olist hứa 24,4 ngày, giao 12,6 ngày, sớm hơn lời hứa 11,8 ngày}
\end{frame}

\begin{frame}{DB2 – Dashboard: trang 2 review và khách quay lại}
\muc{4 · DB2 Giao hàng và trải nghiệm khách}
\begin{itemize}
  \item Cột điểm theo nhóm trễ, donut thanh toán, RFM, cohort
\end{itemize}
\khungchen{CHÈN ẢNH CHỤP DB2 trang 2}

\note{[S101] Có sáu phân khúc RFM trên 94.983 khách. Tiêu đề trước đây: DB2 trang 2: review 4,09, 14,7\% tiêu cực, chỉ 3,1\% khách quay lại}
\end{frame}

\begin{frame}{DB2 – Dashboard: trang 2 cơ cấu thanh toán}
\muc{4 · DB2 Giao hàng và trải nghiệm khách}
\begin{itemize}
  \item 66,9\% giao dịch thẻ trả góp nhiều kỳ
\end{itemize}
\khungchen{CHÈN ẢNH CHỤP DB2 trang 2}

\note{[S102] Boleto chiếm 17,9\%. Tiêu đề trước đây: DB2 trang 2: thẻ tín dụng 78,3\% giá trị thanh toán, trả góp trung bình 3,5 kỳ}
\end{frame}

\begin{frame}{DB2 – Dashboard: trang 3 Key Influencers}
\muc{4 · DB2 Giao hàng và trải nghiệm khách}
\begin{itemize}
  \item Kỳ vọng từ khám phá dữ liệu: giao trễ đứng đầu (Cramér's V 0,403; $\eta$ 0,392)
  \item Mô hình thống kê tự động, không chứng minh nhân quả
\end{itemize}
\khungchen[0.3\textheight]{CHÈN ẢNH CHỤP DB2 trang 3 (Key Influencers)}

\note{[S145] Người 2 dựng trang và ghi các yếu tố đứng đầu. Tiêu đề trước đây: DB2 trang 3: Key Influencers cho biết yếu tố nào đi kèm review tiêu cực}
\end{frame}

\begin{frame}{DB2 – Phát hiện: ba số cần nhớ}
\muc{4 · DB2 Giao hàng và trải nghiệm khách}
\begin{itemize}
  \item Olist hứa 24,4 ngày, giao 12,6 ngày, sớm hơn lời hứa 11,8 ngày
  \item Đúng hạn tụt còn 87,6\% ở Black Friday 11/2017 và 81,0\% ở 03/2018
  \item Review 4,09, 14,7\% tiêu cực, chỉ 3,1\% khách quay lại
\end{itemize}
\note{[P12] Ba phát hiện có số của DB2; số lấy từ các slide Dashboard cùng nhóm và từ phần khám phá dữ liệu.}
\end{frame}

\begin{frame}{DB2 – Quyết định: hành động đề xuất}
\muc{4 · DB2 Giao hàng và trải nghiệm khách}
\begin{itemize}
  \item Giảm trễ ở các tháng tải cao
  \item Ưu tiên xử lý đơn có nguy cơ trễ vì điểm review rơi từ 4,29 xuống 1,70 khi trễ trên 7 ngày
  \item Theo dõi Key Influencers để biết yếu tố đi kèm review tiêu cực
\end{itemize}
\note{[P12] Đề xuất hành động rút từ ba phát hiện; nhóm điều chỉnh theo cách hiểu nghiệp vụ của mình.}
\end{frame}

\begin{frame}{DB2 – DAX chính}
\muc{4 · DB2 Giao hàng và trải nghiệm khách}
\khungchen[0.3\textheight]{CHÈN ẢNH / ĐOẠN DAX: measure chính của DB2 (powerbi/measures.md)}
\note{[P12] Khung chờ ảnh đoạn DAX chính; nguồn powerbi/measures.md (Người phụ trách dashboard điền).}
\end{frame}

\begin{frame}{DB3 – Tình huống: địa lý và phí vận chuyển}
\muc{4 · DB3 Địa lý và phí vận chuyển}
\begin{itemize}
  \item Ai dùng: vận hành và logistics
  \item Câu hỏi: địa lý quyết định trải nghiệm giao hàng và phí vận chuyển ra sao
\end{itemize}
\note{[P12] Ai dùng và câu hỏi của DB3 (mạch Tình huống, Dashboard, Phát hiện, Quyết định).}
\end{frame}

\begin{frame}{DB3 – Dashboard: trang 1 bản đồ và phí}
\muc{4 · DB3 Địa lý và phí vận chuyển}
\begin{itemize}
  \item Bản đồ bong bóng, ma trận phân cấp, GMV khách và seller
\end{itemize}
\khungchen{CHÈN ẢNH CHỤP DB3 trang 1}

\note{[S99] Dùng \texttt{v\_\allowbreak{}geo\_\allowbreak{}state} và \texttt{v\_\allowbreak{}geo\_\allowbreak{}zip} cho bản đồ. Tiêu đề trước đây: DB3 trang 1: Sudeste 65,4\% GMV, Norte và Nordeste chịu phí vận chuyển 22\%}
\end{frame}

\begin{frame}{DB3 – Dashboard: trang 1 drill-down vùng đến thành phố}
\muc{4 · DB3 Địa lý và phí vận chuyển}
\begin{itemize}
  \item Khoảng cách TB 598 km
\end{itemize}
\khungchen{CHÈN ẢNH CHỤP DB3 trang 1}

\note{[S100] Mô hình vai trò địa lý seller cho phép so sánh hai phía. Tiêu đề đầy đủ trước khi rút gọn: DB3 trang 1: 36\% dòng hàng có seller và khách cùng bang; drill-down vùng đến thành phố Tiêu đề trước đây: DB3 trang 1: 36\% dòng hàng seller và khách cùng bang; drill-down tới thành phố}
\end{frame}

\begin{frame}{DB3 – Dashboard: trang 2 drill-through và tooltip}
\muc{4 · DB3 Địa lý và phí vận chuyển}
\begin{itemize}
  \item Chuột phải một bang $\rightarrow$ Drill through $\rightarrow$ Chi tiết bang
  \item Thẻ GMV, Freight ratio, Avg lead days, Late rate; bảng theo thành phố
\end{itemize}
\khungchen[0.3\textheight]{CHÈN ẢNH CHỤP DB3 trang 2 (Chi tiết bang)}

\note{[S146] Người 3 dựng trang drill-through và tooltip. Tiêu đề đầy đủ trước khi rút gọn: DB3 trang 2: drill-through từ một bang sang trang chi tiết bang, kèm tooltip theo bang Tiêu đề trước đây: DB3 trang 2: drill-through từ một bang sang chi tiết bang, kèm tooltip}
\end{frame}

\begin{frame}{DB3 – Phát hiện: ba số cần nhớ}
\muc{4 · DB3 Địa lý và phí vận chuyển}
\begin{itemize}
  \item Sudeste chiếm 65,4\% GMV; Norte và Nordeste chịu phí vận chuyển 22\%
  \item 36\% dòng hàng có seller và khách cùng bang
  \item Khoảng cách trung bình seller $\rightarrow$ khách là 598 km
\end{itemize}
\note{[P12] Ba phát hiện có số của DB3; số lấy từ các slide Dashboard cùng nhóm và từ phần khám phá dữ liệu.}
\end{frame}

\begin{frame}{DB3 – Quyết định: hành động đề xuất}
\muc{4 · DB3 Địa lý và phí vận chuyển}
\begin{itemize}
  \item Ưu tiên cải thiện giao hàng ở Norte và Nordeste
  \item Cân nhắc seller gần khách hơn để giảm phí và lead time
\end{itemize}
\note{[P12] Đề xuất hành động rút từ ba phát hiện; nhóm điều chỉnh theo cách hiểu nghiệp vụ của mình.}
\end{frame}

\begin{frame}{DB3 – DAX chính}
\muc{4 · DB3 Địa lý và phí vận chuyển}
\khungchen[0.3\textheight]{CHÈN ẢNH / ĐOẠN DAX: measure chính của DB3 (powerbi/measures.md)}
\note{[P12] Khung chờ ảnh đoạn DAX chính; nguồn powerbi/measures.md (Người phụ trách dashboard điền).}
\end{frame}

\begin{frame}{DB4 – Tình huống: thu hút và hiệu quả seller}
\muc{4 · DB4 Thu hút và hiệu quả seller}
\begin{itemize}
  \item Ai dùng: quản lý seller và đội sales
  \item Câu hỏi: kênh nào chuyển đổi tốt và seller từ funnel bán ra sao
\end{itemize}
\note{[P12] Ai dùng và câu hỏi của DB4 (mạch Tình huống, Dashboard, Phát hiện, Quyết định).}
\end{frame}

\begin{frame}{DB4 – Dashboard: trang 1 phễu và tỷ lệ chốt}
\muc{4 · DB4 Thu hút và hiệu quả seller}
\begin{itemize}
  \item Phễu MQL $\rightarrow$ deal $\rightarrow$ seller có đơn; thời gian chốt
\end{itemize}
\khungchen{CHÈN ẢNH CHỤP DB4 trang 1 TRANG 1}

\note{[S103] Luôn kèm cỡ mẫu: email chỉ 493 MQL. Tiêu đề trước đây: DB4 trang 1: tỷ lệ chốt theo kênh từ 2,7\% (other) đến 16,7\% (unknown)}
\end{frame}

\begin{frame}{DB4 – Dashboard: trang 2 scorecard seller}
\muc{4 · DB4 Thu hút và hiệu quả seller}
\begin{itemize}
  \item Scorecard seller, GMV/seller theo kênh
\end{itemize}
\khungchen{CHÈN ẢNH CHỤP DB4 trang 2 TRANG 2}

\note{[S104] Email chỉ có 6 seller có đơn nên đánh dấu mẫu nhỏ. Tiêu đề trước đây: DB4 trang 2: seller từ funnel chỉ tạo 4,99\% GMV; kênh nhỏ được cảnh báo mẫu nhỏ}
\end{frame}

\begin{frame}{DB4 – Dashboard: trang 3 chi tiết seller}
\muc{4 · DB4 Thu hút và hiệu quả seller}
\begin{itemize}
  \item Trường drill-through dim\_seller[seller\_id]
  \item Đường GMV theo tháng nhờ snapshot dày đặc, kèm cảnh báo mẫu nhỏ
\end{itemize}
\khungchen[0.3\textheight]{CHÈN ẢNH CHỤP DB4 trang 3 (Chi tiết seller)}

\note{[S147] Người 4 dựng trang chi tiết seller. Tiêu đề trước đây: DB4 trang 3: drill-through sang chi tiết seller cho thấy tháng seller ngừng bán}
\end{frame}

\begin{frame}{DB4 – Phát hiện: ba số cần nhớ}
\muc{4 · DB4 Thu hút và hiệu quả seller}
\begin{itemize}
  \item Tỷ lệ chốt theo kênh từ 2,7\% (other) đến 16,7\% (unknown)
  \item Seller từ funnel chỉ tạo 4,99\% GMV
  \item Kênh email chỉ có 493 MQL và 6 seller có đơn: cỡ mẫu nhỏ
\end{itemize}
\note{[P12] Ba phát hiện có số của DB4; số lấy từ các slide Dashboard cùng nhóm và từ phần khám phá dữ liệu.}
\end{frame}

\begin{frame}{DB4 – Quyết định: hành động đề xuất}
\muc{4 · DB4 Thu hút và hiệu quả seller}
\begin{itemize}
  \item Đánh giá kênh kèm cỡ mẫu, không chỉ theo tỷ lệ
  \item Hỗ trợ 462 deal đã chốt nhưng chưa bán gì
\end{itemize}
\note{[P12] Đề xuất hành động rút từ ba phát hiện; nhóm điều chỉnh theo cách hiểu nghiệp vụ của mình.}
\end{frame}

\begin{frame}{DB4 – DAX chính}
\muc{4 · DB4 Thu hút và hiệu quả seller}
\khungchen[0.3\textheight]{CHÈN ẢNH / ĐOẠN DAX: measure chính của DB4 (powerbi/measures.md)}
\note{[P12] Khung chờ ảnh đoạn DAX chính; nguồn powerbi/measures.md (Người phụ trách dashboard điền).}
\end{frame}

\begin{frame}{Publish và QR: chia sẻ bốn dashboard bằng một liên kết}
\muc{4 · Xây dựng chương trình}
\begin{itemize}
  \item Dữ liệu Olist công khai nên chấp nhận được
\end{itemize}
\khungchen[0.3\textheight]{CHÈN HÌNH: ảnh Publish to web và mã QR của link}

\note{[S106] Nếu tài khoản bị chặn thì dùng PDF xuất từ Power BI. Tiêu đề trước đây: Publish to web chia sẻ cả bốn dashboard bằng một liên kết công khai và mã QR}
\end{frame}

\begin{frame}{Checklist dashboard: dự báo, Key Influencers, OLAP, drill-through}
\muc{4 · Xây dựng chương trình}
\begin{itemize}
  \item Số đơn 12.3K, tiền 13.5M, phần trăm một chữ số thập phân
  \item Màu đỏ chỉ để cảnh báo; theme không còn đỏ trong màu dữ liệu
  \item Dự báo ở DB1 (chuỗi 20 tháng nên chỉ để thấy hướng đi); Key Influencers ở DB2
  \item Ma trận mục checklist theo dashboard ở huong-dan/dashboards/README.md
\end{itemize}

\note{[S143] Checklist rút từ nhận xét mẫu của barem. Tiêu đề đầy đủ trước khi rút gọn: Checklist chung: dự báo ở DB1, Key Influencers ở DB2, OLAP và drill-through ở DB3, DB4 Tiêu đề trước đây: Checklist: dự báo ở DB1, Key Influencers ở DB2, OLAP và drill-through ở DB3, DB4}
\end{frame}
